\documentclass[10pt,twocolumn,a4paper]{article}
\usepackage[pagenumbers]{cvpr}
%% This file contains a number of tweaks that are typically applied to the main document.
%% They are not enabled by default, but can be enabled by uncommenting the relevant lines.


%%
%% Official support for 2-column teaser figure
%%
\usepackage{cuted} %< for strip environment

%%
%% For auto-referencing of figures (see fig/teaser.tex and \label{} therein) 
%%
\usepackage{currfile} %< for \currfilebase macro

%%
%% caption and (global) spacing overrides
%%
\usepackage{caption} %< for "\captionof" commands (teaser)
% \setlength{\abovecaptionskip}{.5em} % space above caption
% \setlength{\belowcaptionskip}{1em} % space below caption

%%
%% Inline annotations; for predefined colors, refer to "dvipsnames" in the xcolor package:
%% https://tinyurl.com/overleaf-colors
%%
\newcommand{\red}[1]{{\color{red}#1}}
\newcommand{\todo}[1]{{\color{red}#1}}
\newcommand{\TODO}[1]{\textbf{\color{red}[TODO: #1]}}
%%
%% Disable for camera ready / submission by uncommenting these lines:
%%
% \renewcommand{\TODO}[1]{}
% \renewcommand{\todo}[1]{#1}

%%
%% Work harder in optimizing text layout. Typically shrinks text by 1/6 of page, enable
%% it at the very end of the writing process, when you are just above the page limit.
%%
% \usepackage{microtype}

%%
%% Fine-tune paragraph spacing.
%%
% \renewcommand{\paragraph}[1]{\vspace{.5em}\noindent\textbf{#1.}}

%%
%% Globally adjusts space between figure and caption.
%%
% \setlength{\abovecaptionskip}{.5em}

%%
%% Allows "the use of \paper to refer to the project name"
%% with automatic management of space at the end of the word.
%%
% \usepackage{xspace}
% \newcommand{\paper}{ProjectName\xspace}

%%
%% Commonly used math definitions.
%%
% \DeclareMathOperator*{\argmin}{arg\,min}
% \DeclareMathOperator*{\argmax}{arg\,max}

%%
%% Tighten underline.
%%
% \usepackage{soul}
% \setuldepth{foobar}

%%
%% Reconfigure paragraph (tweaks spacing)
%%

\usepackage{booktabs,graphicx,xcolor,amsmath}
\definecolor{cvprblue}{rgb}{0.21,0.49,0.74}
\usepackage[breaklinks,colorlinks,allcolors=cvprblue]{hyperref}

% ---------------------------------------------------------------- results plumbing
\def\RES{../results}
\newcommand{\tab}[1]{\IfFileExists{\RES/tables/#1}{\input{\RES/tables/#1}}{\fbox{\texttt{#1} missing: run run\_all.sh}}}
\newcommand{\fig}[2]{\IfFileExists{\RES/figs/#1}{\includegraphics[width=#2]{\RES/figs/#1}}{\fbox{\texttt{#1} missing}}}
\IfFileExists{\RES/tables/numbers.tex}{\input{\RES/tables/numbers.tex}}{}
\IfFileExists{\RES/tables/commit.tex}{\input{\RES/tables/commit.tex}}{\newcommand{\CommitHash}{\textbf{FILL-ME}}}
\newcommand{\RepoURL}{\url{https://github.com/Prashant606468/Machine_Learning_assignment_3}}

\def\confName{CVPR}\def\confYear{2027}
\title{MLPs, CNNs, Autoencoders, SVMs and Clustering on CIFAR10\\
\large Machine Learning, Assignment 3}
\author{Prashant\\ Roll No.\ M24MA2009, IIT Jodhpur\\ {\tt\small m24ma2009@iitj.ac.in}\\
{\small Code: \RepoURL\quad Final commit: \texttt{\CommitHash}}}

\begin{document}
\maketitle

\begin{abstract}
We study why fully connected networks underperform on CIFAR10 and what convolutional and
autoencoder models actually extract from the images. Every comparison fixes its fairness
criterion and its expected outcome before the result is read. Q1 separates the missing spatial
prior of an MLP from a capacity explanation and tracks memorisation of noisy labels through
training. Q2 compares a CNN and an MLP with matched parameter counts under clean, block-shuffled,
pixel-permuted, low-pass and high-pass images. Q3 varies the bottleneck of a convolutional
autoencoder, tests whether a denoising autoencoder restores what a fixed classifier needs, and
compares reconstruction error with linear-probe accuracy. Q4 runs SVMs and clustering on one
frozen autoencoder representation.
At matched parameters the CNN reaches \QcCnnCleanTest\% test accuracy against \QcMlpCleanTest\%
for the MLP; a locality-only change to the MLP's first layer recovers 5.4 of those points, and
pixel permutation removes almost all of the CNN's advantage. An SVM on the frozen
autoencoder codes reaches \QeTest\%, while clustering the same codes barely recovers the classes
(ARI \QeARI).
\end{abstract}

%======================================================================
\section{Common setup}
\label{sec:setup}
\paragraph{Data and split.} The 50{,}000 CIFAR10 training images are split once into a stratified
45{,}000/5{,}000 train/validation split with \texttt{sklearn.model\_selection.train\_test\_split},
\texttt{random\_state}=242009 (the digits of roll number M24MA2009, in order). The index arrays are
saved in \texttt{splits/split\_seed242009.npz} and every script loads that file. Validation is used
for every choice (learning rate, epoch, bottleneck, weight decay, $k$, $C$, kernel, clustering
method). The official test set is evaluated once per final model, after its choices are fixed;
in each lr sweep only the selected run is ever evaluated on test.

\paragraph{Classifier protocol (Q1, Q2).} All MLPs and CNNs share one protocol: per-channel
standardisation with statistics from the (transformed) training split, no data augmentation,
AdamW, one-cycle schedule, 40 epochs, batch 256, weight decay $5\times10^{-4}$, and the learning
rate grid $\{10^{-3},3\times10^{-3}\}$, so each model gets the same tuning budget. The reported
checkpoint is the epoch and learning rate with the highest validation accuracy. Augmentation is
left out on purpose: random crops and flips are themselves a spatial prior, and they would not
commute with the pixel permutation of Q2(b). Train curves use eval-mode accuracy on a fixed
5{,}000-image subset of the training split.

\paragraph{Reuse.} Results are cached by run name, so a model trained once is never retrained.
MLP-base (Q1a) is the MLP of Q2; MLP-base and LocalMLP on pixel-permuted data (Q1a) are reused in
Q2(b); the clean CNN of Q2(a) is the fixed classifier of Q3(b); the selected clean autoencoder of
Q3(a) defines the Q4 representation. Hardware: one NVIDIA RTX 5000 Ada GPU, PyTorch 2.6.0 with
CUDA 12.4; the Q4 SVMs run on the CPU (the full $C\times$kernel sweep took 13.4 minutes in
parallel).

%======================================================================
\section{Q1: Why does an MLP do poorly?}
\subsection{(a) Missing spatial prior vs.\ missing capacity}
\paragraph{Hypotheses.} H1 (chosen): an MLP treats the 3072 inputs as an unordered vector, so
locality and translation structure must be learned from 45k images instead of being built in.
H2 (rival): the MLP simply lacks capacity and underfits.

\paragraph{Design.} Three models on clean data, two of them repeated on pixel-permuted data
(one fixed permutation of the 1024 positions, identical for train, validation and test):
\begin{itemize}\setlength\itemsep{0pt}
\item \textbf{MLP-base}: 3072--512--512--10, BN+ReLU, \QaBaseParams{} parameters.
\item \textbf{MLP-wide}: 3072--2048--512--10, \QaWideParams{} parameters ($4\times$). This is the
capacity probe and deliberately breaks the parameter match.
\item \textbf{LocalMLP}: one linear layer $48\!\to\!48$ shared by all 64 non-overlapping $4\times4$
patches (a $4\times4$, stride-4 convolution), followed by exactly the MLP-base head. It has
\QaLocalParams{} parameters (+0.13\%). The only change is locality and weight sharing in the first
layer.
\end{itemize}
\paragraph{Fairness.} MLP-base and LocalMLP match in parameters (within 0.2\%), depth of the head,
data, epochs and tuning budget. MACs differ slightly and are listed. MLP-wide differs on purpose.

\paragraph{Predictions and decision rule (fixed before training).} Under H2, MLP-base stays far
from 100\% train accuracy and MLP-wide gains clearly on validation. Under H1, MLP-base fits the
training set well while validation stays low, MLP-wide gains little, LocalMLP beats MLP-base on
clean data, and that advantage disappears on permuted pixels, where the patches contain unrelated
pixels. MLP-base should give the same accuracy on clean and permuted data up to seed noise, since
its function class is permutation invariant. We accept H1 over H2 only if all three hold: the
LocalMLP gain on clean data is larger than on permuted data, MLP-base train accuracy is high, and
the MLP-wide gain is small compared with the LocalMLP gain.

\begin{table}[t]\centering\small\resizebox{\linewidth}{!}{\tab{q1a.tex}}
\caption{Q1(a). Train accuracy is eval-mode at the final epoch; validation and test use the
validation-selected checkpoint. Permuted runs reuse the learning rate selected on clean data.}
\label{tab:q1a}\end{table}
\begin{figure}[t]\centering\fig{q1a_curves.pdf}{\linewidth}
\caption{Q1(a) train and validation curves. Dotted: pixel-permuted data.}\label{fig:q1a}\end{figure}

\paragraph{Result.} Test accuracy: MLP-base \QaBaseTest\%, MLP-wide \QaWideTest\%, LocalMLP
\QaLocalTest\%; on permuted data MLP-base \QaBasePermTest\% and LocalMLP \QaLocalPermTest\%.
Final train accuracy of MLP-base: \QaBaseTrainFinal\%.
All three conditions of the decision rule hold. LocalMLP gains $+5.44$ points over MLP-base on
clean data but only $+1.28$ on permuted data. MLP-base fits the training set completely
(\QaBaseTrainFinal\%) while validation stays at \QaBaseVal\%, so it is not underfitting. Quadrupling
the parameters (MLP-wide) buys $+1.90$ points, about a third of what the locality change buys at
equal size. MLP-base gives the same accuracy on clean and permuted pixels ($57.05$ vs.\ $57.09$),
as its permutation invariance requires. We therefore accept H1: the gap comes mainly from the
missing spatial prior. Capacity explains a small part of it.

\subsection{(b) Is structure learned before memorisation?}
\paragraph{Setting.} MLP-wide (7.35M parameters, about 163 per training image), no weight decay,
150 epochs, Adam lr $10^{-3}$. A fixed 40\% of training labels are replaced by a random
\emph{wrong} class. Controls: 0\% noise, and 100\% labels drawn uniformly and independently of
the image (an ``always wrong'' label would leak the true class). Every epoch we measure, on fixed
subsets of the training split, (i) accuracy on clean-label points, (ii) the fraction of
noisy-label points whose \emph{wrong} label is predicted (memorisation), (iii) the fraction of
noisy-label points for which the \emph{true} label is predicted, and (iv) validation accuracy
against true labels.

\paragraph{Prediction.} If useful structure comes first, the early epochs show (iii) and (iv)
rising while (ii) stays near the 10\% chance level; later (ii) rises, (iii) falls and
validation accuracy declines from its peak.

\begin{figure}[t]\centering\fig{q1b_memorisation.pdf}{\linewidth}
\caption{Q1(b). Left: 40\% noise, all four measurements; dotted line at the epoch of best
validation accuracy. Right: validation accuracy for each noise rate and train accuracy on fully
random labels.}\label{fig:q1b}\end{figure}
\begin{table}[t]\centering\small\resizebox{\linewidth}{!}{\tab{q1b.tex}}
\caption{Q1(b). ep$^*$: epoch of best validation accuracy. Mem: fraction of noisy training points
fitted to their wrong label.}\label{tab:q1b}\end{table}

\paragraph{Result.} With 40\% noise the best validation accuracy is \QbNFortyBestVal\% at epoch
\QbNFortyBestEp, when only \QbNFortyMemAtBest\% of noisy points are fitted to their wrong label.
Memorisation passes 50\% at epoch \QbNFortyEpMemFifty{} and ends at \QbNFortyMemFinal\%, while
validation falls to \QbNFortyFinalVal\%. Test accuracy: \QbNFortyTestBest\% (early-stopped) vs.\
\QbNFortyTestFinal\% (final). With random labels the network still reaches high train accuracy and
validation stays at chance, so the capacity to memorise is there.
The plot (Fig.~\ref{fig:q1b}) matches the prediction, with one qualification. Over the first three
epochs the network predicts the \emph{true} label for 44\% of the noisy points and validation
peaks, while the wrong label is fitted at about the 1-in-9 chance rate. Memorisation does not wait
for that peak, though: it climbs steadily from 7\% at epoch 1, so the two phases overlap rather than
being cleanly separated. Between epochs 10 and 30 memorisation rises from 19\% to 89\%, true-label
predictions on noisy points fall to 4\%, and validation drops to 32.5\%. It then partly recovers to
38\% once the training set is interpolated. Early stopping on validation is worth
$+7.7$ test points here. With random labels, train accuracy passes 50\% at epoch
\QbNFullEpMemFifty{} and reaches 100\%, while validation stays at 10--11\%.

%======================================================================
\section{Q2: CNN vs.\ MLP}
\subsection{(a) Fair comparison}
\paragraph{Fairness criterion (fixed before training).} Parameter count within $\pm5\%$ (asserted
in code), identical data and preprocessing, no augmentation, same optimiser, schedule, epochs,
batch size, learning-rate grid and selection rule. Compute cannot be matched at the same time as
parameters: the CNN reuses its weights at every position and does about $100\times$ more MACs.
We report this and treat it as part of what weight sharing means, not as a hidden advantage.
\paragraph{Models.} MLP-base from Q1(a), reused. CNN: three stages of two $3\times3$
conv--BN--ReLU layers (64, 128, 352 channels), $2\times2$ max-pooling after each stage, global
average pooling, linear classifier.
\begin{table}[t]\centering\small\resizebox{\linewidth}{!}{\tab{q2_params.tex}}
\caption{Q2(a) models. MLP-base is reused from Q1(a).}\label{tab:q2p}\end{table}
\paragraph{Proposed reason.} The CNN wins because its filters are local and shared across
positions, so it exploits correlations between neighbouring pixels and reuses one detector
everywhere. Q2(b) tests this.

\subsection{(b) Same pixels, disturbed spatial structure}
Two transforms keep every pixel value and change only positions (Fig.~\ref{fig:q2t}), applied
identically to train, validation and test: \emph{block-shuffle} permutes the sixteen $8\times8$
blocks with one fixed permutation (local structure kept, global layout destroyed);
\emph{pixel-perm} permutes all 1024 positions (no spatial structure left). Both models are
retrained with their clean-selected learning rate.
\paragraph{Expected result (stated before training).} MLP: no change beyond seed noise under
either transform. CNN: a moderate drop under block-shuffle, and under pixel-perm a fall to roughly
the MLP level or below, since its $3\times3$ filters then see unrelated pixels and pooling mixes
them. If the CNN kept most of its advantage under pixel-perm, the explanation of (a) would be
wrong.

\begin{figure}[t]\centering\fig{q2_transforms.pdf}{0.8\linewidth}
\caption{The five conditions on four test images (high-pass shown shifted by 0.5).}\label{fig:q2t}\end{figure}

\subsection{(c) Shape or texture?}
\paragraph{Two explanations.} E1: the CNN relies mostly on low-frequency content (coarse shape,
colour layout). E2: it relies mostly on high-frequency content (edges, local texture).
\paragraph{Experiment.} Low-pass = Gaussian blur, $\sigma=1.5$ px; high-pass = image minus its
low-pass. Both models are retrained and tested on each version. E1 predicts that the low-pass CNN
stays close to the clean CNN and the high-pass CNN loses much more; E2 predicts the reverse. As
an extra test with a distribution shift, the clean CNN is evaluated on low-pass and high-pass test
images (Tab.~\ref{tab:q2x}). Errors are analysed per class and through confusion matrices,
confidence through ECE (15 bins), mean confidence on correct and wrong predictions, and
reliability diagrams.

\begin{table}[t]\centering\small\resizebox{\linewidth}{!}{\tab{q2_main.tex}}
\caption{Q2 test results. Each model is trained and tested under the same condition. conf R/W:
mean max-softmax on right / wrong test predictions.}\label{tab:q2}\end{table}
\begin{table}[t]\centering\small\resizebox{\linewidth}{!}{\tab{q2_cross.tex}}
\caption{Clean-trained models evaluated on transformed test images (shift, not retrained).}\label{tab:q2x}\end{table}
\begin{table}[t]\centering\small\tab{q2_perclass.tex}
\caption{Per-class test accuracy (\%).}\label{tab:q2c}\end{table}
\begin{figure}[t]\centering\fig{q2_results.pdf}{\linewidth}\\[2pt]\fig{q2_confusions.pdf}{\linewidth}
\caption{Top: test accuracy per condition and CNN reliability diagrams. Bottom: CNN confusion
matrices (row-normalised).}\label{fig:q2}\end{figure}

\paragraph{Results.} Clean: CNN \QcCnnCleanTest\% vs.\ MLP \QcMlpCleanTest\%. Block-shuffle: CNN
\QcCnnBlockTest\%, MLP \QcMlpBlockTest\%. Pixel-perm: CNN \QcCnnPermTest\%, MLP \QcMlpPermTest\%.
Low-pass CNN \QcCnnLowTest\% (ECE \QcCnnLowECE), high-pass CNN \QcCnnHighTest\% (ECE
\QcCnnHighECE).
\paragraph{(b) against the expectation.} Every part of the stated expectation holds. The MLP
is unchanged within 0.1 points under both transforms. The CNN loses 11.3 points under
block-shuffle and keeps a 20-point lead, because its $3\times3$ filters still see intact
$8\times8$ neighbourhoods. Under pixel-perm it falls 33.9 points to 54.46\%, 2.6 points
\emph{below} the MLP. Its calibration also degrades, with ECE rising from 0.081 to 0.336. The
advantage in (a) therefore depends almost entirely on pixels keeping their neighbours.

\paragraph{(c) E1 or E2?} The numbers support E2. Retrained on high-pass images the CNN loses only
0.7 points (87.67\%), while on low-pass images it loses 6.6 (81.77\%). The shift test shows the
same thing more sharply. The clean CNN still scores 78.42\% on high-pass test images but only
26.43\% on low-pass ones. Blurring hurts the animal classes most, including those the CNN already
found hardest: dog ($-11$), horse ($-10$), bird ($-9$), deer ($-8$) and cat ($-7$). Vehicles and
ships lose 3--6 points. Wrong predictions stay confident everywhere. Mean confidence on errors is
0.81--0.84 for every CNN, and under shift the clean CNN averages 0.83 confidence on low-pass images
while being right only 26\% of the time (ECE 0.564), and 0.82 on permuted images at 10.7\% accuracy
(ECE 0.713). The MLP shows the opposite frequency preference: it improves slightly on low-pass
images (59.05\%) and collapses on high-pass ones (43.87\%).

\paragraph{Conclusion linking the three experiments.} At equal parameter count the CNN is 31.3
points better than the MLP (a). Experiment (b) shows that this gap depends on spatial arrangement:
keeping local neighbourhoods preserves two-thirds of it, while a full permutation removes all of
it. Experiment (c) shows what those local filters read. The high-frequency band (edges and local
texture) carries nearly all the information the CNN uses, so a $3\times3$ weight-shared detector of
local texture is what the MLP is missing.

%======================================================================
\section{Q3: Autoencoders}
\subsection{(a) Bottleneck capacity}
Encoder: three stride-2 $4\times4$ convolutions (32, 64, 128 channels, BN, ReLU) to
$128\times4\times4$, then a linear map to $d$. The decoder mirrors it with transposed convolutions
and a sigmoid output. Only $d\in\{32,256,1024\}$ changes (compression $96\times$, $12\times$,
$3\times$); the parameter count changes only through the two linear layers next to the
bottleneck. Training: pixel MSE on $[0,1]$ images, Adam lr $10^{-3}$, cosine schedule, 30 epochs,
batch 256, best epoch by validation MSE.
\begin{table}[t]\centering\small\resizebox{\linewidth}{!}{\tab{q3a.tex}}
\caption{Q3(a). Test error is reported only for the model selected on validation.}\label{tab:q3a}\end{table}
\begin{figure}[t]\centering\fig{q3a_curves.pdf}{0.75\linewidth}
\caption{Validation MSE per epoch.}\label{fig:q3a}\end{figure}
\paragraph{Choice.} With reconstruction error as the only criterion, AE-\QdSelD{} has the lowest
validation MSE and is selected; its test MSE is \QdTestMSE{} (PSNR \QdTestPSNR{} dB).
Going from $d=32$ to $d=256$ cuts validation MSE by a factor of 3.4 (PSNR $+5.3$ dB). Going on to
$d=1024$ cuts it by a further 22\% ($+1.1$ dB) at 3.3 times the parameters. All three models were
still improving slightly at epoch 30 (best epoch 29), so the ranking, not the absolute error, is
what we rely on.

\subsection{(b) Denoising}
\paragraph{Corruption and protocol.} Additive Gaussian noise, $\sigma=\QdSigma$ on $[0,1]$
pixels, clipped. Reconstruction can plausibly help because the noise is independent per pixel
while images are spatially smooth. The DAE has the same architecture and bottleneck as the
selected AE; each mini-batch receives fresh noise, the target is the clean image, no labels are
used; the epoch is chosen by MSE on a fixed noisy validation set. The test noise is drawn once
with a fixed seed. The classifier is the clean CNN of Q2(a), frozen. As a non-learned baseline we
add a Gaussian blur of the noisy image.
\begin{table}[t]\centering\small\tab{q3b.tex}
\caption{Q3(b). Fixed CNN on test inputs. Agree: share of predictions equal to the prediction on
the clean image. MSE/PSNR against the clean image.}\label{tab:q3b}\end{table}
\begin{figure}[t]\centering\fig{q3_recons.pdf}{\linewidth}
\caption{Test images: clean, noisy, AE(noisy), DAE(noisy), AE(clean).}\label{fig:q3r}\end{figure}
\paragraph{Decision.} The relevant comparison is DAE(noisy) against the noisy input (does it
help?) and against AE(clean) (the best the bottleneck allows), not against clean images, since
even a clean reconstruction loses high frequencies the CNN uses.
Noise with $\sigma=0.1$ drops the fixed CNN from 88.36\% to 34.76\%. The CNN is very sensitive to
exactly the high-frequency band that the noise corrupts, as Q2(c) showed. The DAE helps, raising
accuracy to 44.44\% ($+9.7$ points) and beating the blur baseline (42.77\%). It stays 6.2 points
short of AE(clean), the best the bottleneck allows. However, the plain AE, never trained on noise,
does slightly better on the same noisy inputs (45.97\%), even though the DAE reaches lower MSE
($0.00356$ vs.\ $0.00362$). Lower reconstruction error therefore did not mean higher accuracy. A
1024-dimensional bottleneck already discards most of the noise, and what limits accuracy is the
high-frequency detail that every reconstruction loses. The blur baseline fits this picture: its MSE
is 22\% higher than the DAE's, but its accuracy is only 1.7 points lower. (The clean-image row is
one image off the 88.34\% of Q2 because inputs are stored in half precision here.)

\subsection{(c) Class information in the codes}
Bottleneck codes of every encoder are frozen and standardised with training statistics. Two
read-outs: multinomial logistic regression (weight decay from
$\{0,10^{-4},10^{-3},10^{-2}\}$ on validation) and cosine $k$-NN ($k\in\{5,20,50\}$ on
validation). An untrained encoder of the same shape is the baseline. The probe trained on clean
codes is also tested on codes of noisy test images.
\begin{table}[t]\centering\small\tab{q3c.tex}
\caption{Q3(c). Test results for frozen encoders.}\label{tab:q3c}\end{table}
\begin{figure}[t]\centering\fig{q3c_mse_vs_probe.pdf}{0.75\linewidth}
\caption{Reconstruction error against linear-probe accuracy (test).}\label{fig:q3c}\end{figure}
\paragraph{Do the two agree?} Only partly. Across the three clean AEs the orderings match: lower
MSE goes with higher probe accuracy (36.06, 44.69, 52.34\%). The DAE breaks the pattern. Its MSE is
15\% worse than AE-1024's, yet its probe is better on clean codes (53.76\%) and on noisy codes
(52.78 vs.\ 49.76\%). It also has the best $k$-NN accuracy (50.96\%). The random encoder is the
clearest warning: with no training at all its 1024-dimensional codes give a 47.87\% probe, above
the trained AE-32 and AE-256. Much of what the probe reads is therefore width plus random features,
and pixel MSE mainly rewards low-level detail. Training on reconstruction adds 4.5 probe points over
random at $d=1024$, and denoising adds another 1.4.

%======================================================================
\section{Q4 (bonus): SVMs and clustering}
\paragraph{Fixed representation.} Codes of AE-\QdSelD{} (Q3a), standardised and reduced by PCA to
128 dimensions (\QePCAVar\% of variance), both fitted once on the 45k training split. All SVMs use
one stratified training subset of \QeSubset{} images, drawn with the split seed before any SVM was
fit and saved in \texttt{splits/}; validation is the full 5k split. $\gamma$ for polynomial
(degree 3, $c_0=1$) and RBF kernels is fixed to $1/(D\,\mathrm{Var}(X))$ on the subset.

\subsection{(a) Effect of $C$ and the kernel}
\begin{table}[t]\centering\small\tab{q4_sweep.tex}
\caption{Q4(a) sweep. $\dagger$: stopped at the iteration cap.}\label{tab:q4a}\end{table}
\begin{figure*}[t]\centering\fig{q4a_sweep.pdf}{0.9\linewidth}
\caption{Validation and train accuracy and support-vector fraction against $C$.}\label{fig:q4a}\end{figure*}
The selected SVM ($\QeKernel$, $C=\QeC$) reaches \QeTest\% test accuracy.
For the RBF and polynomial kernels the sweep has the textbook shape. At $C=0.01$ the margin is
wide, train and validation are close and both low (RBF: 34.4 / 32.0\%), and every training point
is a support vector, which is underfitting. Raising $C$ to 1 lifts validation to its peak (RBF 51.24\%,
poly 50.12\%). Beyond that, train accuracy goes to 99--100\% while validation slips by 0.4--1.5
points, which is mild overfitting. Unlike the textbook case, the support-vector fraction stays at
0.90--0.96 throughout. The classes overlap so heavily in this representation that even a very
tight fit leaves most points on or inside the margin. The linear kernel peaks at the smallest $C$
(41.4\%), about 10 points below the non-linear kernels. Its drop at $C\geq10$ is not overfitting:
those fits hit the iteration cap without converging ($\dagger$), and at $C=100$ train accuracy also
collapses to 19.7\%.

\subsection{(b) Difficult class pairs}
The three pairs with the most test confusions of the multi-class SVM are inspected with the
binary SVM on those two classes (identical to the one-vs-one sub-problem since $\gamma$ is fixed).
\begin{table}[t]\centering\small\tab{q4_pairs.tex}
\caption{Q4(b). MC conf.: multi-class test confusions in both directions. bounded: SVs with
$\alpha=C$. in margin: test points with $|f|<1$. Fisher: squared centroid distance over summed
class variance.}\label{tab:q4b}\end{table}
\begin{figure*}[t]\centering\fig{q4b_pairs.pdf}{0.85\linewidth}
\caption{Per pair: the most violating training SVs ($y f$ smallest) and the eight test images
closest to the boundary, with decision scores $f$.}\label{fig:q4b}\end{figure*}
\begin{figure}[t]\centering\fig{q4b_margins.pdf}{\linewidth}
\caption{Test decision scores for each difficult pair.}\label{fig:q4m}\end{figure}
\textbf{Cat--dog} (405 confusions) is the hardest pair, with a Fisher ratio of 0.008, five times
below the all-pair mean of 0.042. 92\% of its training points are support vectors, 89\% of test
points fall inside the margin, and the binary SVM still errs on 35\%. The most violating SVs are
pale, fluffy close-ups: white cats on grass on one side, and frontal faces of long-haired small dogs
that look much like cats on the other. In other words, colour and texture are shared and the face
is the only cue. \textbf{Automobile--truck} (310) separates best of the three (Fisher 0.022, binary
error 23.5\%). Its boundary cases are pickups, vans and white box trucks photographed side-on
against plain backgrounds, with the same boxy silhouette and pale body. \textbf{Bird--deer} (293)
is almost as inseparable as cat--dog (Fisher 0.009, binary error 30.8\%). Both classes appear as
small brown animals in green or brown vegetation, and several near-boundary test images are
labelled frog, cat or dog by the multi-class SVM, so the code mostly encodes the scene. In all
three pairs 91--94\% of the multi-class errors lie inside the binary margin. The errors are
therefore points of genuine overlap, not confident mistakes far from the boundary.

\subsection{(c) Clustering}
$k$-means and diagonal GMMs with $k\in\{5,10,15,20,30\}$ are fit on the training representations
without labels; method and $k$ are chosen by silhouette on validation assignments, also without
labels. The selected model assigns the test representations; test labels are used only afterwards.
\begin{table}[t]\centering\small\tab{q4_cluster_sel.tex}
\caption{Q4(c) clustering selection.}\label{tab:q4c}\end{table}
\begin{table}[t]\centering\small\tab{q4_contingency.tex}
\caption{Test cluster-to-class contingency table.}\label{tab:q4ct}\end{table}
\begin{figure}[t]\centering\fig{q4c_contingency.pdf}{0.85\linewidth}
\caption{P(cluster $\mid$ class) on test.}\label{fig:q4c}\end{figure}
\begin{table}[t]\centering\small\tab{q4_disagree.tex}\\[3pt]\tab{q4_class.tex}
\caption{Classification vs.\ clustering. Overlap $=\sum_c\min(P(c|i),P(c|j))$; rank gap $>0$:
clustering merges a pair the SVM separates.}\label{tab:q4d}\end{table}
Selected: \QeClusterMethod, $k=\QeK$; test ARI \QeARI, NMI \QeNMI.
Silhouettes are close to zero for every $k$ and negative for all GMMs, so the codes have no
well-separated cluster structure. $k$-means with $k=5$ is chosen because it is the only setting
with a positive score. The clusters follow scene type rather than class. One cluster collects
trucks, ships, automobiles and horses (Tab.~\ref{tab:q4ct}, cluster 1); another collects most
birds, deer and frogs (cluster 4). The clearest disagreement is \textbf{automobile--dog}. Their
cluster distributions overlap by 0.824, so clustering hardly tells them apart, yet the SVM confuses
only 1.4\% of their test images. The autoencoder was trained to reproduce pixels, so the largest
directions of variance in its codes are global colour, brightness and background, and $k$-means
splits along those. The SVM uses labels to find the lower-variance directions that carry the class.
The reverse also occurs. Airplane and ship are confused by the SVM 12\% of the time but are
\emph{better} separated by the clusters than average (rank gap $-0.52$), because both share
blue-sky or blue-sea backgrounds, which the SVM cannot use to tell them apart.

%======================================================================
\section{Reproducibility}
\texttt{run\_all.sh} runs every experiment and rebuilds this PDF; all hyperparameters live in
\texttt{configs/config.yaml}; every number in the text is generated into
\texttt{results/tables/numbers.tex}; logs (per-epoch histories, sweeps, test logits) are in
\texttt{results/logs}. Seeds are fixed, but cuDNN autotuning makes GPU runs non-bitwise
deterministic, so retraining can move accuracies by a few tenths of a point.

\end{document}
